\documentclass[12pt,a4paper]{article}
\usepackage[utf8]{inputenc}
\usepackage[T1]{fontenc}
\usepackage[brazil]{babel}
\usepackage{indentfirst}
\usepackage[margin=2.5cm]{geometry}

\title{Resenha Cr\'itica: Artigo 7 ---\\
Domain-Driven Design Reference}
\author{Gustavo Pessoa Firmino Duarte}
\date{12 de Abril de 2026}

\begin{document}
\maketitle

\section{Introdu\c{c}\~ao}

O \textit{Domain-Driven Design Reference} (2015) de Eric Evans \'e um guia de
refer\^encia conciso que sintetiza os padr\~oes e conceitos fundamentais do
\textit{Domain-Driven Design} (DDD), abordagem que Evans introduziu em seu livro
hom\^onimo de 2003. O documento n\~ao pretende substituir o livro original, mas
oferecer uma refer\^encia r\'apida e precisa para praticantes que j\'a conhecem
o paradigma. Seu valor est\'a em tornar expl\'icitos os padr\~oes que o DDD defende:
desde a constru\c{c}\~ao de uma \textit{linguagem ub\'iqua} compartilhada entre
desenvolvedores e especialistas do dom\'inio, passando pelo isolamento do modelo
em \textit{contextos delimitados} (\textit{bounded contexts}), at\'e os blocos
construtivos t\'aticos como entidades, objetos de valor, agrega\c{c}\~oes,
reposit\'orios e servi\c{c}os de dom\'inio.

\section{Linguagem Ub\'iqua, Contextos Delimitados e Blocos T\'aticos}

O conceito mais transformador do DDD \'e a \textit{linguagem ub\'iqua}
(\textit{ubiquitous language}): uma linguagem comum, rigorosa e sem ambiguidade,
constru\'ida colaborativamente por desenvolvedores e especialistas do neg\'ocio,
que deve permear o c\'odigo, a documenta\c{c}\~ao e as discuss\~oes da equipe.
Evans argumenta que a separa\c{c}\~ao entre a ``linguagem dos neg\'ocios'' e a
``linguagem do c\'odigo'' \'e uma fonte constante de erros e mal-entendidos;
quando o c\'odigo usa os mesmos termos que os especialistas usam no dia a dia, a
descoberta de inconsist\^encias no modelo de dom\'inio torna-se mais f\'acil.

Os \textit{bounded contexts} definem as fronteiras dentro das quais um modelo de
dom\'inio \'e v\'alido e consistente. Em sistemas grandes, diferentes sub-dom\'inios
possuem modelos distintos e frequentemente conflitantes para o mesmo conceito ---
um ``cliente'' em um contexto de faturamento pode ter atributos completamente
diferentes de um ``cliente'' em um contexto de suporte. O DDD abraça essa
multiplicidade, desde que as fronteiras sejam expl\'icitas e as tradu\c{c}\~oes
entre contextos sejam gerenciadas atrav\'es de padr\~oes como \textit{Anti-Corruption
Layer}, \textit{Shared Kernel} ou \textit{Open Host Service}.

No n\'ivel t\'atico, Evans distingue \textbf{entidades} (objetos com identidade
ao longo do tempo) de \textbf{objetos de valor} (\textit{value objects}, imut\'aveis
e definidos apenas por seus atributos). As \textbf{agrega\c{c}\~oes} (\textit{aggregates})
s\~ao clusters de entidades e objetos de valor tratados como uma unidade transacional,
acessados exclusivamente atrav\'es de uma raiz de agrega\c{c}\~ao. \textbf{Reposit\'orios}
abstraem o mecanismo de persist\^encia, e \textbf{eventos de dom\'inio}
(\textit{domain events}) capturam acontecimentos significativos no dom\'inio,
viabilizando integra\c{c}\~ao assíncrona entre contextos.

\section{Conex\~ao com a Pr\'atica Profissional}

Ao desenvolver o MVP de e-commerce com Node.js, Prisma ORM e React, implementei
estruturas que, em retrospecto, guardam semelhan\c{c}a com os padr\~oes t\'aticos
do DDD --- ainda que sem o rigor conceitual que Evans propugna. Os modelos do
Prisma atuavam como entidades, e usei servi\c{c}os para encapsular a l\'ogica de
neg\'ocio. Por\'em, a aus\^encia de agrega\c{c}\~oes bem definidas resultou em
opera\c{c}\~oes que modificavam m\'ultiplas entidades relacionadas fora de
transa\c{c}\~oes at\^omicas, criando janelas de inconsist\^encia de dados.

Na ArcelorMittal, a diversidade de sistemas internacionais que suportei ilustrou
claramente a problem\'atica dos contextos delimitados: o mesmo conceito de ``ordem''
(\textit{order}) tinha semânticas distintas no sistema de produ\c{c}\~ao brasileiro,
no sistema log\'istico europeu e no sistema financeiro global. A aus\^encia de
\textit{bounded contexts} expl\'icitos tornava as integra\c{c}\~oes entre esses
sistemas fr\'ageis, com regras de tradu\c{c}\~ao espalhadas e freq\"uentemente
duplicadas.

A linguagem ub\'iqua \'e um aspecto que percebo como especialmente dif\'icil de
cultivar: requer disciplina cont\'inua para que termos do dom\'inio sejam usados
com consist\^encia no c\'odigo, nos tickets e nas discuss\~oes --- uma pr\'atica
que raramente vi implementada de forma sistem\'atica em projetos reais.

\section{Conclus\~ao}

O \textit{Domain-Driven Design Reference} de Evans \'e um compendio essencial para
qualquer desenvolvedor que trabalhe com sistemas complexos. Seus padr\~oes ---
especialmente bounded contexts, agrega\c{c}\~oes e linguagem ub\'iqua --- oferecem
ferramentas concretas para domar a complexidade acidental que frequentemente contamina
o modelo de dom\'inio em sistemas de longa vida. O DDD n\~ao \'e uma bala de prata:
demanda investimento em compreens\~ao do neg\'ocio e disciplina de design. Mas para
sistemas com dom\'inios ricos e equipes colaborativas, os benef\'icios em clareza,
manutenibilidade e agilidade diante de mudan\c{c}as de neg\'ocio s\~ao substanciais.

\end{document}
